% =====================================================================
%  chapters/0-coverpage.tex
%  Khóa luận tốt nghiệp: "Phương pháp tạo dữ liệu kiểm thử thủ công
%  cho đối tượng atomic trong dự án C++"
%  Sinh viên: Lưu Quang Tùng - MSSV 23020150
%  Cán bộ hướng dẫn: PGS.TS. Phạm Ngọc Hùng
%
%  Gồm: 3 trang bìa (bìa có logo, bìa có cán bộ hướng dẫn, bìa tiếng Anh),
%       Lời cam đoan, Lời cảm ơn, Tóm tắt (tiếng Việt) và Abstract (tiếng Anh).
%
%  LƯU Ý: dòng "HÀ NỘI - <năm>" được đẩy xuống sát đáy khung nhờ lệnh \vfill
%  đặt ngay trước nó. Nếu bạn thêm/bớt nội dung phía trên, dòng này vẫn ở đáy.
% =====================================================================

%-------TITLE PAGE------%
\begin{titlepage}
	\center{}
	\begin{tikzpicture}[overlay,remember picture]
		\draw [line width=3pt,rounded corners=0pt,]
		($ (current page.north west) + (25mm,-25mm) $) rectangle ($ (current page.south east) + (-15mm,25mm) $);
		\draw [line width=1pt,rounded corners=0pt]
		($ (current page.north west) + (26.5mm,-26.5mm) $) rectangle ($ (current page.south east) + (-16.5mm,26.5mm) $);
	\end{tikzpicture}
	
	{\large \bfseries ĐẠI HỌC QUỐC GIA HÀ NỘI\\ TRƯỜNG ĐẠI HỌC CÔNG NGHỆ}\\[1cm]
	\includegraphics[width=0.25\linewidth]{images/Logo_UET.png}\\[1cm]
	
	{\Large \bfseries Lưu Quang Tùng}\\[1.5cm]
	{\LARGE \bfseries ECLIPSE CDT VÀ CLANG: NGHIÊN CỨU THỰC NGHIỆM ĐỊNH HƯỚNG THIẾT KẾ LỚP PHÂN TÍCH MÃ NGUỒN TRONG CÔNG CỤ KIỂM THỬ C/C++}\\[0.2cm]
	% {\LARGE \bfseries CHO ĐỐI TƯỢNG ATOMIC TRONG DỰ ÁN C++}\\[0.2cm]
	
	\vspace{1.5cm}
	{\large \bfseries KHÓA LUẬN TỐT NGHIỆP ĐẠI HỌC HỆ CHÍNH QUY}\\
	{\large \bfseries Ngành: Công nghệ thông tin}
	
	\vfill
	{\large \bfseries HÀ NỘI \@- \the\year}\\
	\vspace{0.5cm}
\end{titlepage}

%-------TITLE PAGE (kèm cán bộ hướng dẫn)------%
\begin{titlepage}
	\center{}
	\begin{tikzpicture}[overlay,remember picture]
		\draw [line width=3pt,rounded corners=0pt,]
		($ (current page.north west) + (25mm,-25mm) $) rectangle ($ (current page.south east) + (-15mm,25mm) $);
		\draw [line width=1pt,rounded corners=0pt]
		($ (current page.north west) + (26.5mm,-26.5mm) $) rectangle ($ (current page.south east) + (-16.5mm,26.5mm) $);
	\end{tikzpicture}
	
	{\large \bfseries ĐẠI HỌC QUỐC GIA HÀ NỘI\\ TRƯỜNG ĐẠI HỌC CÔNG NGHỆ}\\[2cm]
	
	{\Large \bfseries Lưu Quang Tùng}\\[2cm]
	{\LARGE \bfseries ECLIPSE CDT VÀ CLANG: NGHIÊN CỨU THỰC NGHIỆM ĐỊNH HƯỚNG THIẾT KẾ LỚP PHÂN TÍCH MÃ NGUỒN TRONG CÔNG CỤ KIỂM THỬ C/C++}\\[0.2cm]
	% {\LARGE \bfseries CHO ĐỐI TƯỢNG ATOMIC TRONG DỰ ÁN C++}\\[0.2cm]
	
	\vspace{1.5cm}
	{\large \bfseries KHÓA LUẬN TỐT NGHIỆP ĐẠI HỌC HỆ CHÍNH QUY}\\
	{\large \bfseries Ngành: Công nghệ thông tin}\\[1cm]
	\begin{flushleft}
		{\large \bfseries Cán bộ hướng dẫn: PGS.TS. Phạm Ngọc Hùng}
	\end{flushleft}
	
	\vfill
	{\large \bfseries HÀ NỘI \@- \the\year}\\
	\vspace{0.5cm}
\end{titlepage}

%-------TITLE PAGE (tiếng Anh)------%
\begin{titlepage}
	\center{}
	\begin{tikzpicture}[overlay,remember picture]
		\draw [line width=3pt,rounded corners=0pt,]
		($ (current page.north west) + (25mm,-25mm) $) rectangle ($ (current page.south east) + (-15mm,25mm) $);
		\draw [line width=1pt,rounded corners=0pt]
		($ (current page.north west) + (26.5mm,-26.5mm) $) rectangle ($ (current page.south east) + (-16.5mm,26.5mm) $);
	\end{tikzpicture}
	
	{\large \bfseries VIETNAM NATIONAL UNIVERSITY, HANOI\\ UNIVERSITY OF ENGINEERING AND TECHNOLOGY}\\[2cm]
	
	{\Large \bfseries Luu Quang Tung}\\[2cm]
	{\LARGE \bfseries ECLIPSE CDT AND CLANG: AN EMPIRICAL STUDY ON THE DESIGN OF THE SOURCE CODE ANALYSIS LAYER IN C/C++ TESTING TOOLS}\\[0.2cm]
	% {\LARGE \bfseries FOR ATOMIC OBJECTS IN C++ PROJECTS}\\[0.2cm]
	
	\vspace{1.5cm}
	{\large \bfseries BACHELOR'S THESIS}\\
	{\large \bfseries Major: Information Technology}\\[1cm]
	\begin{flushleft}
		{\large \bfseries Supervisor: Assoc. Prof., PhD Pham Ngoc Hung}
	\end{flushleft}
	
	\vfill
	{\large \bfseries Hanoi \@- \the\year}\\
	\vspace{0.5cm}
\end{titlepage}

\changefontsizes[16pt]{13pt}

\addtocontents{toc}{\vspace{-1cm}}
\addcontentsline{toc}{chapter}{Lời cam đoan}
\begin{center}
    \textbf{LỜI CAM ĐOAN}
\end{center}

Em xin cam đoan: Khóa luận tốt nghiệp ``Eclipse CDT và Clang: Nghiên cứu thực nghiệm định hướng thiết kế lớp phân tích mã nguồn trong công cụ kiểm thử C/C++''
là của em. Toàn bộ những tri thức, nội dung, kết quả của người khác có đề cập trong khóa luận đều được trích dẫn nguồn gốc
cụ thể. Khóa luận và mã nguồn là công trình nghiên cứu và phát triển của cá nhân em.

Nếu vi phạm những điều đã cam đoan, em xin chấp nhận tất cả những truy cứu về trách nhiệm theo quy định của Trường Đại
học Công nghệ \@- ĐHQGHN.\@

\begin{flushright}
    \begin{varwidth}{\linewidth}\centering Hà Nội, ngày \ldots{} tháng \ldots{} năm \the\year\\
        Sinh viên\\[2cm]
        Lưu Quang Tùng
    \end{varwidth}
\end{flushright}

\newpage

\addcontentsline{toc}{chapter}{Lời cảm ơn}
\begin{center}
    \textbf{LỜI CẢM ƠN}
\end{center}

Lời đầu tiên cho phép em bày tỏ lòng biết ơn chân thành đến Khoa Công nghệ Thông tin \@- Trường Đại học Công nghệ, ĐHQGHN,
nơi em đã có cơ hội tiếp xúc với những kiến thức mới mẻ và được học hỏi từ những người giáo viên xuất sắc.

Em xin gửi lời cảm ơn chân thành nhất đến PGS.TS. Phạm Ngọc Hùng, những người đã tận tình định
hướng, góp ý và động viên em trong suốt quá trình thực hiện khóa luận. Sự hướng dẫn tận tâm và hỗ trợ sát sao của các
thầy là nguồn động lực lớn giúp em vượt qua những khó khăn trong quá trình nghiên cứu.

Cuối cùng, em xin gửi lời cảm ơn tới các thành viên của Phòng thí nghiệm Đảm bảo chất lượng phần mềm (Khoa Công nghệ thông
tin, Trường Đại học Công nghệ, ĐHQGHN) đã tư vấn, ủng hộ và động viên em trong quá trình hoàn thành khóa luận.

\newpage
\addcontentsline{toc}{chapter}{Tóm tắt}
\begin{center}
    \textbf{TÓM TẮT}
\end{center}
\changefontsizes[16pt]{12pt}

\textit{\textbf{Tóm tắt:}} Trong phát triển phần mềm hiện đại, kiểm thử đơn vị đóng vai trò quan trọng trong việc đảm bảo
chất lượng và độ tin cậy của hệ thống. Công cụ AKAUTAUTO, được phát triển bởi Phòng thí nghiệm Đảm bảo chất lượng phần mềm
(Trường Đại học Công nghệ, ĐHQGHN) phối hợp cùng FPT-GAM, sử dụng kỹ thuật thực thi tượng trưng động để hỗ trợ kiểm thử tự
động cho các chương trình C/C++. Tuy nhiên, hệ thống vẫn gặp hạn chế khi xử lý các chương trình sử dụng đối tượng
\texttt{atomic}, do chưa hỗ trợ đầy đủ việc mô hình hóa và tương tác với kiểu dữ liệu này trong quá trình phân tích và
kiểm thử.

Nghiên cứu này đề xuất mở rộng AKAUTAUTO nhằm hỗ trợ các đối tượng \texttt{atomic} trong quy trình kiểm thử. Giải pháp
tập trung vào hai hướng chính:

\begin{itemize}
    \item \textbf{Mở rộng mô hình biểu diễn dữ liệu nguyên tử:} bổ sung cơ chế nhận diện và xây dựng cấu trúc dữ liệu
    tương ứng cho các biến \texttt{atomic} trong quá trình phân tích chương trình, cho phép hệ thống hiển thị và quản lý
    các đối tượng nguyên tử tương tự các kiểu dữ liệu đã được hỗ trợ trước đó.
    \item \textbf{Hỗ trợ thực thi và đánh giá kiểm thử đối với các đối tượng nguyên tử:} bổ sung khả năng xử lý các thao
    tác phổ biến như \texttt{load}, \texttt{store}, \texttt{exchange} và \texttt{compare\_exchange}, cho phép người dùng
    thiết lập dữ liệu kiểm thử, thực thi chương trình và theo dõi kết quả độ phủ mã nguồn đối với các nhánh liên quan đến
    biến \texttt{atomic}.
\end{itemize}

Kết quả thực nghiệm cho thấy giải pháp đề xuất giúp AKAUTAUTO hỗ trợ hiệu quả hơn việc kiểm thử các chương trình sử dụng
đối tượng nguyên tử. Hệ thống có thể nhận diện và biểu diễn chính xác các biến \texttt{atomic}, đồng thời hỗ trợ người
dùng xây dựng và thực thi các ca kiểm thử để đánh giá độ phủ mã nguồn. Cải tiến này góp phần mở rộng phạm vi áp dụng của
AKAUTAUTO đối với các hệ thống nhúng và ứng dụng đa luồng sử dụng cơ chế đồng bộ hóa dựa trên \texttt{atomic}.

\vspace{-0.5cm}
\begin{flushleft}
  \textit{\textbf{Từ khóa: }AKAUTAUTO, kiểm thử đơn vị, thực thi tượng trưng động, \texttt{atomic}, độ phủ mã nguồn, C/C++.}
\end{flushleft}

\newpage
\addcontentsline{toc}{chapter}{Abstract}
\begin{center}
    \textbf{ABSTRACT}
\end{center}
\changefontsizes[16pt]{12pt}

\textit{\textbf{Abstract:}} In modern software development, unit testing plays a key role in ensuring the quality and
reliability of software systems. The AKAUTAUTO tool, developed by the Software Quality Assurance Laboratory (University of
Engineering and Technology, VNU Hanoi) in cooperation with FPT-GAM, applies dynamic symbolic execution to support automated
testing of C/C++ programs. However, the tool still has limitations when handling programs that use \texttt{atomic} objects,
because the modelling of and interaction with this data type are not fully supported during program analysis and testing.

This study proposes an extension of AKAUTAUTO to support \texttt{atomic} objects in the testing workflow. The solution
focuses on two main directions: (i) extending the representation model for atomic data --- adding a mechanism that
recognises \texttt{atomic} variables and builds the corresponding data structure during program analysis, so that the
system can display and manage atomic objects similarly to previously supported data types; and (ii) supporting test
execution and evaluation for atomic objects --- adding handling for common operations such as \texttt{load}, \texttt{store},
\texttt{exchange} and \texttt{compare\_exchange}, allowing users to configure test data, execute the program and monitor
code coverage for branches related to \texttt{atomic} variables.

Experimental results show that the proposed solution enables AKAUTAUTO to support the testing of programs that use atomic
objects more effectively. The system recognises and represents \texttt{atomic} variables correctly, and helps users build
and execute test cases for evaluating code coverage. This improvement extends the applicability of AKAUTAUTO to embedded
systems and multithreaded applications that rely on atomic-based synchronisation.

\vspace{-0.5cm}
\begin{flushleft}
  \textit{\textbf{Keywords: }AKAUTAUTO, unit testing, dynamic symbolic execution, \texttt{atomic}, code coverage, C/C++.}
\end{flushleft}
\changefontsizes[16pt]{13pt}
